\subsubsection{Isothermal Viscous\slash No-Slip Wall}
\label{SEC:BC:Wall:Isothermal}

The isothermal viscous\slash no-slip wall boundary condition enforces a
prescribed wall velocity and temperature at the wall boundary surface.
%
The isothermal no-slip wall boundary condition is very similar to the adiabatic
no-slip wall in \sref{SEC:BC:Wall:Adiabatic}, with the only difference between
them being how the wall temperature and temperature gradients are set.
%
Therefore, the remainder of this section only provides details for what is done
differently from the adiabatic wall boundary condition.
%
Particularly, all of the details in \sref{SEC:BC:Wall:Adiabatic} on how the
prescribed wall velocity is determined are exactly the same for isothermal
walls.

%%%%%%%%%%%%%%%%%%%%%%%%%%%%%%%%%%%%%%%%%%%%%%%%%%%%%%%%%%%%%%%%%%%%%%%%%%%%%%%%

The backbone of the isothermal wall boundary condition is prescribing a fixed
wall temperature for a solid boundary surface.
%
The prescribed wall temperature for the boundary condition, $T_\subBC$, is
specified using the input
%\BCinptext{Wall\_Temp}.
\BCinpBC[nmlBC::BCInput/WallTemp]{Wall\_Temp}.
%
For the rest of the thermodynamic gas properties, pressure is taken from the
interior, and then the ideal gas law is used to compute the density at the wall,
resulting in the following right state variables:
%
\begin{alignat*}{1}
  T_\subR & = T_\subBC \\
  p_\subR & = p_\subL \\
  \rho_\subR & = \frac{p_\subR}{\rgas T_\subR}
\end{alignat*}

%%%%%%%%%%%%%%%%%%%%%%%%%%%%%%%%%%%%%%%%%%%%%%%%%%%%%%%%%%%%%%%%%%%%%%%%%%%%%%%%

Finally, unlike for adiabatic walls, the isothermal wall temperature gradient is
just copied from the interior,
%
\begin{equation*}
  \nabla T_\subR = \nabla T_\subL
\end{equation*}
%
Since this is also done for the gradients of all of the other variables, the
right gradient states for the isothermal wall are just the standard
\quoterm{flow-through} gradient boundary condition from
\eref{EQ:BC:FlowThroughGradient} used by most boundary conditions.

\devNotes{If default flow conditions are found for an isothermal wall, it will
use the \inptext{ovar\_mod} variable \inptext{iso\_wall\_temp} to set the wall
temperature to 300 K. This probably needs to be changed to an error+abort. This
dev-note is also in \sref{SEC:BC:Specify:FullState}.}

%###############################################################################

